\documentclass[11pt,a4paper]{article}
\usepackage[T1]{fontenc}
\usepackage{textcomp}
\usepackage[margin=25mm]{geometry}
\usepackage[hidelinks]{hyperref}
\setlength{\parindent}{0pt}
\setlength{\parskip}{6pt}
\setlength{\emergencystretch}{2em}

\title{Baseline RAG and Cross-Encoder Reranking}
\author{Aryan\\Roll No. 2025201017}
\date{}

\begin{document}
\maketitle

\section{Document and retrieval setup}

The source document is Leder Games\textquotesingle{} \emph{The Law of Root},
dated October 13, 2025. A rulebook suited this experiment because a general
rule can be qualified by a restriction several pages later. The PDF has 32
pages. Printed pages 2 to 29 form the corpus; the cover, index, credits and
repeated quick reference are excluded. The notebook verifies the PDF hash
and reads the text in column order.

Both pipelines share 317 overlapping chunks with source-page metadata.
The largest chunk in this run has 230 tokens, within the embedding
model\textquotesingle{}s 256-token limit. The final word limit is 130, with
an overlap target of 32 words.

For baseline retrieval, \texttt{sentence-transformers/all-MiniLM-L6-v2}
embeds each chunk independently. FAISS \texttt{IndexFlatIP} stores the
normalized 384-dimensional vectors. Each query is embedded once.
Inner product on normalized vectors gives cosine similarity, which is
used to retrieve ten chunks. The first three are passed to the LLM for
Answer 1.

The cross-encoder, \texttt{cross-encoder/ms-marco-MiniLM-L-6-v2}, processes
the query and each candidate chunk together. Joint self-attention lets it
compare tokens across the two texts before producing a score. It scores
the same ten candidates, sorts them, and passes the new top three to the
LLM for Answer 2.

The answer model is Groq\textquotesingle{}s \texttt{qwen/qwen3.8-27b}.
Both pipelines use the same prompt, temperature 0.2, reasoning effort
\texttt{none}, and a 1,024-token output limit. The prompt asks for answers
based only on the excerpts, including restrictions and chunk/page
citations. It also asks the model to say when the evidence is insufficient.

\section{Where reranking helped}

Most of the first 15 questions already had enough evidence in the baseline
top three. Several chunks moved up, but those movements did not establish
an answer improvement.

The additional search tested 20 questions across eight rule families.
It screened for a source-rule passage missing from the baseline top three
and promoted into the reranked top three. Three selected questions received
answer comparisons. The source-rule anchors were used for screening only;
they were never added to either answer prompt.

\subsection*{Boat Builders and hirelings}

For the question about a hireling using its
controller\textquotesingle{}s Boat Builders effect, the baseline retrieves
Riverfolk Company rules about river movement. The subject is similar, but
those passages say nothing about a hireling inheriting its
controller\textquotesingle{}s effects. Answer 1 says there is not enough
information to answer.

Reranking moves \texttt{chunk\_0292} from rank 5 to rank 1. On page 26,
rule H.3.5 says hirelings cannot use their
controller\textquotesingle{}s abilities or crafted persistent effects.
It even names Boat Builders as an example. Answer 2 gives the supported
answer: no.

This case shows how passages about boats and river movement can crowd out
a restriction on who may use an effect.

\subsection*{Exposure at the end of a turn}

The second question asks whether a player can expose a Corvid plot after
starting their final draw-and-discard step. The baseline finds draw/discard
instructions and Corvid overview material. Answer 1 cannot determine
whether Exposure is allowed at that point.

The missing restriction is in \texttt{chunk\_0154}, on page 15. It moves
from rank 8 to rank 1. Rule 13.2.4 explicitly says that Exposure cannot be
used once the player has begun the last step of their turn. Answer 2
answers no and cites this rule. Its discussion of the neighboring
draw/discard instructions is less direct, but the Exposure restriction
itself resolves the question.

Here, retrieving the turn sequence was not enough. The answer depended
on a timing condition in a different passage.

\subsection*{A third example: coalition with a Hostile faction}

For the question about joining a Hostile faction and moving its marker to
Indifferent, \texttt{chunk\_0092} moves from rank 4 to rank 1. The baseline
has the general coalition rules but cannot answer the relationship part.
The reranked answer explains that forming a coalition moves the marker to
Indifferent. The rule is on page 10, within a chunk spanning pages 9 to 10.
Normal coalition conditions still apply, including the four-player minimum.

\begin{center}
\small
\begin{tabular}{p{0.40\textwidth}p{0.24\textwidth}rr}
\hline
Question & Relevant chunk & Before & After \\
\hline
Hireling using Boat Builders & \texttt{chunk\_0292} & 5 & 1 \\
Exposure after the final step begins & \texttt{chunk\_0154} & 8 & 1 \\
Coalition with a Hostile faction & \texttt{chunk\_0092} & 4 & 1 \\
\hline
\end{tabular}
\end{center}

In these cases the baseline abstains rather than inventing a rule.
Reranking provides the evidence needed to answer.

\section{Retrieval time}

The notebook ran on Kaggle. After a warm-up, it measured each original
query five times, giving 75 samples. The three selected stress questions
also have five samples each. Baseline time covers query embedding and
FAISS search. Reranked time adds the cost of scoring ten pairs.
Model loading, indexing, answer generation and API retry waits are excluded.
The supplied files do not record the exact Kaggle accelerator.

\begin{center}
\small
\begin{tabular}{p{0.37\textwidth}rrr}
\hline
Measurement & Baseline (ms) & Cross-encoder (ms) & Total (ms) \\
\hline
Original queries, pooled median & 6.26 & 21.10 & 27.50 \\
Boat Builders, median of 5 & 7.14 & 21.29 & 28.43 \\
Late Exposure, median of 5 & 6.56 & 20.69 & 27.23 \\
Hostile coalition, median of 5 & 6.73 & 22.54 & 30.06 \\
\hline
\end{tabular}
\end{center}

The ratio of the original queries\textquotesingle{} median reranked time
to their median baseline time is about 4.39. Most of the additional time
comes from the cross-encoder. Each column reports its own median, so the
component medians need not sum exactly to the median total. These are
measurements of this run; the hardware details are insufficient for a
broader speed comparison.

\section{Why not score every chunk with a cross-encoder?}

The bi-encoder computes passage embeddings once and reuses them for later
queries. The current FAISS flat index performs an exact vector search.
An approximate index could be used for a larger collection.

A cross-encoder must process every passage together with the current query.
It cannot reuse an independent passage embedding in the same way.
For this corpus, scoring everything would require 317 pairs per query
instead of ten, about 32 times as many pair evaluations. Batching and
hardware affect the runtime, so that pair count does not imply an exact
32-fold latency increase. The cost grows with the collection size.
Candidate retrieval keeps this expensive scoring step limited to ten pairs.

\section{What still failed}

The original Hostile-coalition question missed the relevant chunk in its
initial top ten. Both answers lacked the rule. A differently worded stress
question retrieved it and allowed reranking to promote it. This is a
limit of the two-stage setup: the cross-encoder can only reorder what it
receives. Query wording affected which passages reached that stage.

The Winter-map case exposed another problem. The baseline applies the
general rule that rivers do not divide forests. Reranking promotes the
specific exception from rank 5 to rank 2, but the selected chunk begins
at M.2.1 and leaves out its parent Winter Map heading. The LLM treats
the exception as an optional variant and still gives a misleading final
answer. Keeping section headings with their child rules is a possible
improvement. It was not tested here.

All 36 answers finished with \texttt{stop} and no recorded API error.
That confirms completion. Factual accuracy still requires checking the
answers against the source. The successful cases above were checked
against the PDF and their supplied contexts.

The questions were selected through exploratory screening, so these
results do not give an overall accuracy rate on unseen questions.

\section*{Code and supporting files}

Repository: \url{https://github.com/a4Aaaryan/rag-baseline-vs-reranking}

Rulebook source: \url{https://ledergames.com/pages/resources}

The notebook and dependency list are in the repository.
\texttt{answers\_completed.csv} contains the original comparisons;
\texttt{stress\_answers.csv} contains the three additional pairs.
The evidence files retain the retrieved passages and their ranks.
The timing CSVs contain individual measurements. API keys and the
downloaded rulebook are excluded from Git.

\end{document}
